\textbf{Despeje del punto neutro.} La pendiente directa es $C_{m\alpha}=a_w(h-h_{ac})-K_t(h_t-h)+C_{m\alpha,f}$. Agrupar términos en $h$ deja $(a_w+K_t)h-a_wh_{ac}-K_th_t+C_{m\alpha,f}$. Igualar a cero produce la expresión de $h_n$; conservar $K_t$ en el denominador es necesario porque la cola también sustenta.\begin{enumerate}

\item La pendiente efectiva de cola incluye área, presión dinámica relativa y estela.

\[\fitmath{\mathrm{Kt}=0.94\cdot 0.21\cdot 4\cdot \left(1-0.32\right)\approx 0.536928\,}\]

\item Sumar las contribuciones a la sustentación total.

\[\fitmath{\mathrm{atot}=5.2+\mathrm{Kt}}\]

\[\fitmath{\mathrm{atot}=5.2+(0.536928)\approx 5.736928\,}\]

\item Imponer pendiente de momento nula y despejar la posición de CG.

\[\fitmath{\mathrm{hn}=\frac{5.2\cdot 0.25+\mathrm{Kt}\cdot 3.4-0.08}{\mathrm{atot}}}\]

\[\fitmath{\mathrm{hn}=\frac{5.2\cdot 0.25+(0.536928)\cdot 3.4-0.08}{(5.736928)}\approx 0.5308686\,}\]

\item Restar el CG al punto neutro.

\[\fitmath{\mathrm{SM}=\mathrm{hn}-0.38}\]

\[\fitmath{\mathrm{SM}=(0.5308686461)-0.38\approx 0.1508686\,}\]

\item Comprobar signo de la pendiente restauradora.

\[\fitmath{\mathrm{Cm}=-\mathrm{atot}\cdot \mathrm{SM}}\]

\[\fitmath{\mathrm{Cm}=-(5.736928)\cdot (0.1508686461)\approx -0.8655226\,\mathrm{1/rad}}\]

\end{enumerate}
